\chapter{生产消费、位点与重复投递}
本章先区分四种容易混淆的进度。日志末端 LEO 是某分区下一条可追加记录的 offset；复制时，HW（高水位）是下一条 offset 边界，只有小于 HW 的记录处于复制可见范围。在目前的单 broker 路径中 HW 与 LEO 相同，复制章节才会出现两者不同。consumer 的 position 是客户端内存中“下次从哪里拉”的 offset；committed next offset 是按 group 与分区保存、重启后可恢复的进度。poll 只推进内存 position，commit 才写入进度日志；二者都不删除消息。producer receipt 是 broker 对某次追加返回的分区及半开 offset 区间 \pathcode{[firstOffset,nextOffset)}，不等于消费提交。本章的直接 \method{RpcClient} 由客户端接管并在 \method{close} 时关闭；传入的共享 \method{MetadataRouter} 则由创建者负责关闭。

\lessonsection{13}{批量生产者}
\goals{实现按分区有序的批量缓冲、达到阈值时同步发送，并在显式 flush 时交还完整 receipt。}{理解 producer、分区路由、batch、确认与请求结果未知分别是什么。}
\background{producer 是把记录写入 broker 的客户端；每条记录先按 topic（记录分类）和 partition（各自有序的日志）定位。batch 是暂存在客户端、准备一次发送的一组记录；按分区分别缓冲，所以一个分区达到阈值时不会顺便发送其他分区的记录。receipt 是 broker 回复的实际 offset 区间；ACK 是 producer 请求的确认模式。本步的单机 backend 忽略 ACK；非负的 ProduceRequest.timeoutMillis 也不触发等待。producer 同步等待 broker 回复，实际 TCP 等待上限由单独配置的 \method{RpcClient} 控制；这不表示多副本确认，副本 ACK 含义见 \stepref{23}{生产确认}。收到错误或 timeout 并不能证明没有追加：请求可能已到 broker 并落盘，只是回复丢失。课程的普通 \method{SimpleProducer} 不自动重试，也不使用后台定时器。}
\subsubsection*{开始前先读}
先理解 \stepref{3}{分区与分区内 offset}，再复习 \stepref{9}{key 如何选择分区}、\stepref{10}{metadata 如何列出分区}及 \stepref{12}{broker 如何处理 PRODUCE}。特别注意 \method{Partitioner.choose} 是第 9 步学生已完成的代码，并非本步提供的答案；本步当前使用直接 \method{RpcClient} 路径，\method{MetadataRouter} 客户端路由留到第 27 步。
\providededit{\method{RpcClient}、消息/记录值类型和 broker metadata/produce 协议已提供；复用第 9 步实现的 \method{Partitioner}。}{\pathcode{exercises/src/main/java/io/simplekafka/client/SimpleProducer.java}：\method{void send(String topic, byte[] key, byte[] value, long timestamp)} 与 \method{List<ProduceReceipt> flush()}。只实现这两个练习方法，不改 \method{Partitioner}、协议或测试。}
\subsubsection*{参数、结果与状态}
\begin{itemize}
\item 构造参数 \pathcode{batchRecords} 必须大于 0，否则 \method{INVALID_REQUEST}；传入的 client/partitioner 为 null 抛 \pathcode{NullPointerException}。默认 ACK 为 \pathcode{LEADER}，请求字段 \pathcode{timeoutMillis}=5,000 ms。已有 \method{setAcks(Acks,long)} 修改后续请求的 ACK/timeout 字段，timeout 必须非负，且只能在没有待发 batch 时更改；null acks 抛 \pathcode{NullPointerException}。
\item \method{send} 的 topic 必须是课程有效标识；key 可为 null；value 必须非 null；timestamp 是非负的 long，课程不要求把它解释成真实墙钟。成功调用不返回 receipt，而是把记录加入对应 \pathcode{TopicPartition} 的 FIFO batch；达到阈值即同步发送该分区。
\item \method{flush()} 发送所有非空分区 batch，返回自上次 flush 以来积累的不可变 receipt 列表，并清空这份 receipt 累积；没有待发记录时不发空请求，返回空列表。
\item 发送前的 metadata/参数错误没有 broker 追加。进入 produce RPC/回复路径后的错误可能已有追加；producer 保存首次错误并进入结果未知状态，后续 \method{send}、\method{flush}、\method{setAcks} 均报 \method{REQUEST_TIMEOUT}，只能关闭后换新实例。
\end{itemize}
\lessoncontract{相同分区按 send 顺序追加；不同分区各自成 batch。自动发送的 receipt 也必须留到下一次显式 flush，receipt 使用左闭右开区间。null key 由分区器的实例级轮转处理，非 null key 走第 9 步路由。}
\subsubsection*{必须满足的边界}
\begin{itemize}
\item null value、负 timestamp、无效 topic 在进入 produce 之前报 \method{INVALID_REQUEST}；未知 topic 报 \method{UNKNOWN_TOPIC_OR_PARTITION}。这些检查不会追加，也不会让 producer 进入结果未知状态。
\item 超长普通记录（总长超过 1,048,580 bytes；null key 时 value 上限为 1,048,548 bytes）报 \method{INVALID_REQUEST}，且 broker 未收到 produce frame、没有追加。若 batch 尚未达到阈值，\method{send} 可能先正常返回；本地编码器要到自动发送或 \method{flush()} 才发现超长。此时 producer 会把这次 produce 路径异常记为结果未知，后续操作报 \method{REQUEST_TIMEOUT}。这是课程客户端的保守状态行为，不代表 broker 可能已经追加，也不得重发该 batch。
\item 首次发送 batch 的同步路径发生任何异常时，当前调用传播原错误码并锁存 producer 的 unknown-outcome 状态；之后 \method{send}、\method{flush}、\method{setAcks} 均报 \method{REQUEST_TIMEOUT}，不得重发。实际追加是否发生取决于异常位置：网络/回复丢失时可能已经追加；本地编码错误能确定未发送，但本课程实现仍会锁存该状态。
\item 正在缓冲记录时调用 \method{setAcks} 报 \method{INVALID_REQUEST}；负 timeout 也报 \method{INVALID_REQUEST}。关闭后调用 producer 操作报 \pathcode{IllegalStateException}。
\end{itemize}
\expected{手算：broker 的 \pathcode{orders-0} 初始 LEO=0，构造 \pathcode{batchRecords=2} 的 producer。send 的两条同 key 记录依次缓冲，第二条触发发送，broker 追加 offset 0、1，区间为 \pathcode{[0,2)}；第三条缓冲后尚不可读。下一次 \method{flush()} 发送第三条并返回两张 receipt：\pathcode{[0,2)} 与 \pathcode{[2,3)}；再 flush 返回空。前提是 topic 只有该目标分区或 key 已知路由到它。\\
\method{Step13Test.flushesAtTheBatchThresholdAndReturnsAutomaticAndExplicitRanges} 检查阈值与只交还一次的 receipt；\method{automaticFlushForOnePartitionLeavesAnotherPartitionsPendingBatchInvisible} 检查分区 batch 独立；\method{rejectsInvalidSendsWithoutPersistingAnyRecords} 和 \method{doesNotRetryAProduceWhoseAppendSucceededBeforeItsResponseTimedOut} 检查参数拒绝及追加后丢回复不重试。}
\lessonhints{先分清请求发出、broker 追加、客户端收到确认和调用者取得 receipt 的时刻；它们可能不是同一时刻。}{用 \pathcode{TopicPartition} 作 batch 边界，分别手算同分区达到阈值和不同分区尚未达到阈值。}{用测试中的“追加成功后延迟/丢弃回复”情形检查：客户端是否保留首次错误，并且是否把旧 batch 再发一次。}
\lessonquestions{为什么 timeout 不能证明记录没有写入？}{若自动发送已成功，为什么 receipt 还要由下一次 flush 返回？}
\paragraph{自检解释。}请求回复可能在追加之后丢失，所以客户端无法判断结果；receipt 是调用者获知具体 offset 区间的依据，丢掉它会使已追加记录变得不可追踪。
\pitfalls{把“收到错误”理解成“肯定未写入”并重试会产生重复。真实 Kafka 还有异步发送与更多确认选项；本课程的 \method{SimpleProducer} 是同步、显式 flush、无自动重试的简化客户端。}
\lessoncommand{13}

\lessonsection{14}{拉取消费者}
\goals{构建单线程、单轮 pull consumer，显式管理分配分区和内存 position。}{理解“下一条 offset”、seek 重放、空 fetch 与跨分区共享预算。}
\background{consumer 是主动向 broker 拉取记录的客户端；读取不会从日志删除数据。position 是某分区下次 fetch 的 offset，不是已经持久化的消费进度。手动 assign 建立客户端负责读取的分区集合；seek 只改变本地读取位置。poll 发起一轮有上限的 fetch，而不是持续等待。}
\subsubsection*{开始前先读}
复习 \stepref{3}{分区和局部 offset}、\stepref{12}{broker 的 FETCH 请求}。第 13 步只是课程顺序上的前一步，本步不调用 producer，也不读取 OffsetStore；本步先从 offset 0 建立纯内存 position，第 15 步才加入恢复与提交。
\providededit{broker 的 \pathcode{FETCH} 路由、\method{RpcClient}、\method{LogRecord} 与 \method{TopicPartition} 类型已提供。}{\pathcode{exercises/src/main/java/io/simplekafka/client/SimpleConsumer.java}：\method{void assign(Collection<TopicPartition> partitions)}、\method{void seek(TopicPartition tp, long offset)}、\method{Map<TopicPartition,List<LogRecord>> poll(int maxRecords, int maxBytes)}。只实现三者；只读 \method{position(TopicPartition)} 已提供。}
\subsubsection*{参数、结果与状态}
\begin{itemize}
\item \method{assign} 用新集合替换旧 assignment：去重并按 topic/partition 排序；保留分区沿用 position，撤销分区移除，新分区从 0 开始。null 集合抛 \pathcode{NullPointerException}，含 null 元素报 \method{INVALID_REQUEST}，拒绝时旧 assignment 不变。
\item \method{seek(tp,offset)} 只接受已分配 tp 与非负 next offset；null tp 抛 \pathcode{NullPointerException}，未分配报 \method{NOT_ASSIGNED}，负数报 \method{OFFSET_OUT_OF_RANGE}，拒绝时 position 不变。seek 不提交进度。
\item \method{poll} 的 \pathcode{maxRecords} 与 \pathcode{maxBytes} 都必须大于 0，否则 \method{INVALID_REQUEST} 且不推进 position。二者是所有已分配分区共享的总预算；扫描顺序为 topic、partition 升序。空 assignment 不发 fetch，空分区不消耗预算也不出现在返回 map 中；成功返回的 map 不可变。consumer 关闭后调用这些操作报 \pathcode{IllegalStateException}，未分配 partition 的 \method{position(tp)} 报 \method{NOT_ASSIGNED}。
\item 本步普通记录计费为 \pathcode{32 + keyBytes + valueBytes}；不能拆分记录。每分区成功回复须从请求 position 起连续且不超剩余记录/字节预算，否则该分区报 \method{INVALID_REQUEST} 且不前移。若 seek 到已删除的起点之前或当前 LEO 之后，broker 报 \method{OFFSET_OUT_OF_RANGE}，失败分区不推进。服务端其他错误原码传播；先前已成功分区的 position 可保留推进，失败分区保持原值。
\end{itemize}
\lessoncontract{成功返回记录的分区 position 推进到最后一条 offset 加一；未返回的分区、空 fetch 与失败分区不推进。poll 不自动 commit，也不删除日志。过大的正 maxBytes 会由 broker 限制到保守单帧预算；下一轮仍从首个未返回 offset 继续。}
\expected{假设 \pathcode{orders-0}、\pathcode{orders-1} 各有两条普通记录，每条 key=null、value 为空、因此计费 32 bytes；二者 position 初始均为 0。assign 两分区后 poll(4,64) 按排序先读 \pathcode{orders-0} 的 offset 0、1，返回只含该分区的 map，position 变为 (2,0)，预算用尽。把两者 seek 回 0 后 poll(4,63)，第一分区只能返回 offset 0，剩余 31 bytes 放不下下一条，最终 position 为 (1,0)。若之后的脚本化测试 broker 对第二分区返回 \method{UNKNOWN_TOPIC_OR_PARTITION}（模拟该请求失败），第一次 position 可为 1、失败分区仍为 0。\\
\method{Step14Test.pollsInPartitionOrderWithinRoundBudgetsAndPreservesManualPositions} 检查排序、预算、seek 与 assign；\method{pollContinuesPastEmptyPartitionsWithoutSpendingTheSharedBudget} 检查空分区；\method{leavesFailedPartitionPositionUnchangedForFetchErrorsAndMalformedSuccessReplies} 检查失败、断续 offset 和超预算响应。}
\lessonhints{分别记录 assignment、每分区内存 position 和持久 committed offset；本步只改前两者。}{每返回一条记录就从同一个 poll 剩余预算扣除完整编码字节与一条记录，不要给每个分区重置预算。}{先让一个较早排序的分区成功、再让后一个分区失败，逐项检查两个 position，而不要假设整轮原子回滚。}
\lessonquestions{poll 读到 offset 0、1 并把 position 设为 2 后，重建 consumer 为什么仍可能从 0 读？}{为什么后一个分区 fetch 失败时不能把它的 position 也推进？}
\paragraph{自检解释。}position 只存在旧客户端内存中，重建后若没有 resume/commit 就从新 assignment 的默认 0 开始；失败分区没有被成功读取，推进它会跳过数据，而先前成功分区的读取进度可单独保留。
\pitfalls{读取成功不等于持久提交；seek 到日志末端得到空结果，不是让 position 自动增加的理由。真实 Kafka 有持续 poll、心跳及更复杂的分配生命周期，本步只实现单线程、有限轮次。}
\lessoncommand{14}

\lessonsection{15}{消费位点持久化}
\goals{把各 group/分区的下一条 offset 写入本地持久日志，并由新 consumer 恢复。}{区分消费位点、OffsetStore 自己的物理日志 offset，以及 group key。}
\background{consumer position 是客户端内存的读取游标；commit 把“下一条要读取的 offset”保存为消费位点。OffsetStore 自己也用追加日志，但那条日志的物理 offset 只编号存储记录，payload 中的 \pathcode{nextOffset} 才是消费进度。键为 \pathcode{(group,TopicPartition)}，不同 group 或分区互不覆盖；读取和提交都不删除业务记录，也不改变 HW。}
\subsubsection*{开始前先读}
需要理解 \stepref{1}{日志记录与恢复}、\stepref{12}{broker 请求路由}和 \stepref{14}{position、assign、poll}。手动 consumer 用 null group token；本步还没有第 20 步的成员围栏。
\providededit{Offset entry 的 \method{MessageCodec} 编解码、底层 \method{PartitionLog}、记录恢复和 \pathcode{COMMIT_OFFSET/FETCH_OFFSET} 请求类型已提供。}{\pathcode{exercises/src/main/java/io/simplekafka/group/OffsetStore.java}：\method{void commit(OffsetKey key, long nextOffset)}、\method{OptionalLong fetch(OffsetKey key)}；\pathcode{exercises/src/main/java/io/simplekafka/client/SimpleConsumer.java}：\method{void resume(Collection<TopicPartition> partitions)}、\method{void commitSync()}；并在 \pathcode{exercises/src/main/java/io/simplekafka/broker/BrokerHandler.java} 的 \method{handle(Messages.Request)} 增补 \pathcode{COMMIT_OFFSET/FETCH_OFFSET} 分支。}
\subsubsection*{参数、结果与状态}
\begin{itemize}
\item \pathcode{OffsetKey} 标识 group 与一个 topic-partition；\method{commit} 的 nextOffset 必须非负，允许显式回退。commit 的 null key、负数，以及 fetch 的 null key 均报 \method{INVALID_REQUEST}；已关闭 store 或存储故障为 \method{STORAGE_ERROR}。\method{fetch} 返回最新 \pathcode{OptionalLong}；没有提交是 empty，不是 offset 0 的已提交记录。
\item 每次 commit 以 null record key、timestamp=0、编码后的 offset entry 追加并 force；重开时扫描同一 key 的最后一条。物理日志 offset 与 payload 的 nextOffset 没有相等要求。完整损坏记录报 \method{CORRUPT_RECORD} 且不改文件；合法但不完整的尾记录可按底层恢复规则截断。
\item \method{resume} 对分区去重、排序并替换 assignment；每个 position 设成已提交值，缺失时为 0，不把已提交值夹到当前 LEO。位点查询失败会传播，但 assignment 已替换且不会回滚：新分区仍为 0，保留分区保留原 position。 \method{commitSync} 为每个已分配分区各发一次提交，包括 position=0；空 assignment 不发送提交请求。poll 本身不提交，多个分区提交不是一个原子事务。
\item 协议成功响应：\pathcode{CommitOffsetRequest(OffsetKey,nextOffset,null)} 得 \pathcode{EmptyBody}；\pathcode{FetchOffsetRequest(OffsetKey)} 得带 OptionalLong 的 \pathcode{OffsetBody}。负提交值经 broker 也报 \method{INVALID_REQUEST}。
\end{itemize}
\lessoncontract{最后一次追加的同 key 值为当前提交，允许回退；group 和 topic-partition 任一不同都隔离。resume 只恢复持久消费进度，不修复越界提交，也不会自动把 consumer position 写回 store。}
\expected{初始 store 无 \pathcode{workers/orders-0} 值。commit(3) 后关闭并重开，fetch 得 3；再 commit(1) 并重开，fetch 得 1，展示 latest-wins 且可回退。另一 group 的相同分区仍为空。新 consumer.resume 后 position=1；若提交值设成 99，即使分区 LEO=3，resume 仍为 99，随后 poll 才报 \method{OFFSET_OUT_OF_RANGE} 且不改该值。\\
\method{Step15Test.offsetStorePersistsNextOffsetsPerGroupAndPartitionAndAllowsRewind} 检查隔离、回退和重开；\method{rawOffsetRequestsReturnExactBodiesAndIsolateGroupTopicAndPartitionKeys} 检查协议体；\method{resumeCommitsAllPositionsAndDoesNotClampPastLogEnd} 检查 resume；\method{offsetStoreTruncatesAnIncompleteTailAndRejectsCompleteBadCrcWithoutMutation} 检查恢复和完整损坏。}
\lessonhints{先把消费位点看成“下一条业务记录的编号”，不要用 OffsetStore 日志的 offset 代替它。}{同一 OffsetKey 可以有多条历史提交；以追加顺序最后一条为准，不按 map 顺序猜测。}{分开验证两件事：store 关闭重开后值是否仍在，以及全新 consumer.resume 后 position 是否从该值开始。}
\lessonquestions{处理完 offset 2 后为什么提交 nextOffset=3，而不是 2？}{如果 nextOffset=99 而 LEO=3，resume 是否应自动改成 3？}
\paragraph{自检解释。}nextOffset 表示下次要处理的位置，所以 3 才跳过已处理的 0--2；store 只忠实恢复提交值，越界由后续 fetch 明确报错，自动夹取会悄悄改写调用者的进度。
\pitfalls{提交 last processed offset 会让它在恢复后重读；把 committed offset 与复制可见边界混为一谈也不对。真实 Kafka 把 offset 放在复制/压缩的内部主题；本课程是本地单 coordinator 追加日志，不提供高可用。}
\lessoncommand{15}

\lessonsection{16}{至少一次处理}
\goals{安全安排 poll、业务副作用和 offset commit 的先后，并能手算故障后的重放。}{理解至少一次、内存 position 与持久提交的区别，以及跨请求非原子边界。}
\background{业务副作用是处理记录时对外部状态的改变，例如写本地账本或发送通知。课程的 \method{ProcessingLoop} 一次调用先 poll，再按 TopicPartition 升序、同分区 offset 升序调用 processor；所有回调成功后才调用一次 \method{commitSync}。副作用与 offset 日志不是同一事务：若处理已成功而提交失败，恢复后必须重读，可能再次产生副作用。至少一次的含义是宁可重复已成功工作，也不把未成功处理的记录静默跳过。}
\subsubsection*{开始前先读}
需要 \stepref{14}{poll 与 position}、\stepref{15}{持久 next offset 与 commitSync}。本步提供回调接口，不要求外部副作用与 offset store 共享事务。
\providededit{\method{RecordProcessor.process(TopicPartition,LogRecord)} 回调接口与单轮 poll/commit 生命周期已提供。}{\pathcode{exercises/src/main/java/io/simplekafka/client/ProcessingLoop.java}：\method{int runOnce(int maxRecords, int maxBytes) throws Exception}。只实现单轮控制；不改 consumer、store 或回调接口。}
\subsubsection*{参数、结果与状态}
\begin{itemize}
\item maxRecords、maxBytes 是本轮 poll 跨分区的正总预算；非正值由 poll 报 \method{INVALID_REQUEST}。成功返回本轮回调并提交的记录数，空 poll 仍调用 commitSync 并返回 0。
\item callback 输入为分区和一条 \method{LogRecord}，不返回值，可抛任意 \pathcode{Exception}。回调异常原样传播，本轮不提交任何分区；poll 可能已经推进 consumer 的内存 position。
\item poll、callback 或 commit 出现异常后，此 ProcessingLoop 变为 failed；再次 runOnce 报 \pathcode{IllegalStateException}。关闭旧 consumer，创建新 consumer 和 loop，并 resume 已持久提交位置。副作用已经写出但 offset 尚未提交时，副作用可能重做。
\item 所有回调成功后 \method{commitSync} 对每个已分配分区分别提交；不是跨分区事务，某个 commit 失败时较早提交可能已成功。失败原异常/协议错误向调用者传播，方法不返回成功计数。
\end{itemize}
\lessoncontract{回调只按稳定分区、offset 顺序执行；只有所有被本轮 poll 返回的记录都处理成功，才开始提交。一个 callback 失败时不提交任一分区。任何本轮失败后不得复用 loop 或把已推进的内存 position 当作 durable progress。}
\expected{初始 \pathcode{orders-0}、\pathcode{orders-1} 各有 offset 0、1，group \pathcode{workers} 没有提交，poll(4,8192) 返回四条。回调先产生 effects \pathcode{p0:0,p0:1}，再在 \pathcode{p1:0} 抛错；p0/p1 的 committed offset 都仍缺失，尽管旧 consumer 内存 position 已到 2。重建并 resume 后从 0/0 重读四条，之前 p0 两个副作用会再次发生。另一种故障是所有回调完成后 commit 请求失败：副作用已发生而提交未完整确认，恢复仍可能重放；若有多个分区提交，部分 key 可能已写入。\\
\method{Step16Test.laterPartitionFailureLeavesBothOffsetsUncommittedAndReplaysEarlierSideEffects} 覆盖跨分区回调失败；\method{commitFailureAfterSideEffectsReplaysTheSameRecordsAtLeastOnce} 覆盖提交故障；\method{boundedRoundsProcessInTopicPartitionOffsetOrderAndCommitFinalPositions} 覆盖排序、预算、空轮；\method{failedPollProcessesNothingAndResumeRetriesThePreviouslyFetchedRecord} 覆盖 poll 失败。}
\lessonhints{把 poll、每次副作用、每次 offset 请求、崩溃/重建按时间顺序写出。}{先按排序规则列出本轮回调，再判断失败落在 callback 还是 commit；这两者对副作用的影响不同。}{观察 Step16Test 中每个失败后 store 的值、consumer 的 position 和 loop 是否还能调用；分别记录它们，不能用一个“进度”概括。}
\lessonquestions{如果 p0 的回调成功、p1 的回调失败，为什么 p0 的成功副作用也会重放？}{commitSync 对三个分区是不是全成或全败？}
\paragraph{自检解释。}契约等所有回调完成才开始提交，因此 callback 失败时没有任何分区被提交，恢复从旧的 durable offsets 重新读取；commitSync 是多个请求，不是跨分区事务，故较早分区可能已提交。
\pitfalls{吞掉 callback 异常后继续 commit 会丢失失败记录；继续使用失败 loop 则可能跳过未提交数据。至少一次不等于 exactly-once；应用需让副作用可去重，第 31 步才实现本地持久业务去重。}
\lessoncommand{16}
